% !TEX TS-program = xelatex
% !TEX encoding = UTF-8 Unicode

\documentclass{resume}
\usepackage{zh_CN-Adobefonts_external}
\usepackage{linespacing_fix}
\usepackage{xcolor}

\hypersetup{
  colorlinks=true,
  urlcolor=black,
  linkcolor=black
}

\begin{document}
\pagenumbering{gobble}

\name{Junhao Zhao \heitiT{\textbf{赵俊豪}}}

\basicInfo{
  \email{nemozjh@hotmail.com}
  \textperiodcentered\
  \github[nemozjh]{https://github.com/nemozjh}
  \textperiodcentered\
  \homepage[ORCID: 0000-0002-3271-6455]{https://orcid.org/0000-0002-3271-6455}
}
\basicInfo{
  \homepage[nemozjh.github.io]{https://nemozjh.github.io}
}

\section{\faUser\ Research Profile}

\begin{onehalfspacing}
I am an incoming Ph.D. student at \textbf{Xi'an Jiaotong-Liverpool
University} (September 2026), supervised by
\href{https://lijieyao.com}{Dr. Lijie Yao}. My research lies at the
intersection of \textbf{visualization}, \textbf{human-computer interaction},
and \textbf{AI agents}. I am particularly interested in
visualization in motion, interactive authoring tools, and reliable AI-assisted
data communication. Previously, I studied artificial intelligence and 3D
scene perception with
\href{https://asc.xmu.edu.cn/t/wenchengluChengluWen}{Prof. Chenglu Wen}
at Xiamen University.
\end{onehalfspacing}

\section{\faGraduationCap\ Education}

\datedsubsection{\textbf{Xi'an Jiaotong-Liverpool University}, Suzhou, China}
{Incoming, Sep 2026}
\textit{Ph.D. Student}. Advisor:
\href{https://lijieyao.com}{Dr. Lijie Yao}.

\datedsubsection{\textbf{Xiamen University}, Xiamen, China}
{Sep 2020 -- Jun 2023}
\textit{M.Sc. in Artificial Intelligence}. Advisor:
\href{https://asc.xmu.edu.cn/t/wenchengluChengluWen}{Prof. Chenglu Wen}.

\datedsubsection{\textbf{Nanjing Audit University}, Nanjing, China}
{Sep 2014 -- Jun 2018}
\textit{B.Sc. in Information Management and Information Systems}.

\section{\faFlask\ Research Experience}

\datedsubsection{\textbf{Research Assistant, Xi'an Jiaotong-Liverpool University}}
{Apr 2025 -- Present}
\role{Supervisor: \href{https://lijieyao.com}{Dr. Lijie Yao}}{}
\begin{itemize}
  \item Designed and implemented \textbf{SwimComposer}, an interactive
  authoring environment for coordinating data layers, views, transitions, and
  temporal composition in swimming-race videos; the work was accepted by IEEE
  TVCG for presentation at IEEE VIS 2026.
  \item Developed a multimodal data-preparation pipeline that aligns video
  tracking, commentary-derived insights, athlete information, and official
  records into visualization-ready data.
  \item Developed structured AI agents that produce executable and verifiable
  outputs for reliable visualization report generation, published at the
  IEEE VIS 2025 VISxGenAI workshop.
\end{itemize}

\datedsubsection{\textbf{Weakly Supervised Learning for Dynamic 3D Scene Perception}}
{Jan 2022 -- Jun 2023}
\role{NSFC General Program, Grant No. 62171393}{}
\begin{itemize}
  \item Designed weakly supervised learning strategies for sequential point
  clouds and evaluation methods for improving pseudo-label quality through
  multi-perspective temporal analysis.
  \item The work resulted in a first-authored IEEE TGRS paper and a Chinese
  national invention patent application.
\end{itemize}

\datedsubsection{\textbf{Multi-platform Multimodal Point-cloud Data Processing}}
{May 2022 -- Dec 2022}
\role{National Key R\&D Program Young Scientists Project,
Grant No. 2021YFF0700188}{}
\begin{itemize}
  \item Supported multimodal data preprocessing, heterogeneous-platform
  integration, and software development for scalable point-cloud management.
\end{itemize}

\newpage
\section{\faBook\ Publications}

\begin{enumerate}[label={[\arabic*]}]
  \item \textbf{Junhao Zhao}, Romain Vuillemot, Petra Isenberg, and Lijie Yao.
  ``Authoring Narrative Visualization in Motion: Visual Storytelling in
  Swimming Videos.'' \textit{IEEE Transactions on Visualization and Computer
  Graphics}. Accepted; to be presented at \textbf{IEEE VIS 2026}.

  \item \textbf{Junhao Zhao} and Lijie Yao.
  ``Structured AI Agents for Reliable Visualization Report Generation.''
  \textit{IEEE VIS 2025 Workshop on GenAI, Agents, and the Future of VIS
  (VISxGenAI)}.

  \item \textbf{Junhao Zhao}, Weijie Huang, Hai Wu, Chenglu Wen, Bo Yang,
  Yulan Guo, and Cheng Wang.
  ``SemanticFlow: Semantic Segmentation of Sequential LiDAR Point Clouds from
  Sparse Frame Annotations.'' \textit{IEEE Transactions on Geoscience and
  Remote Sensing}, vol. 61, pp. 1--11, 2023.
\end{enumerate}

\section{\faBalanceScale\ Patent}

\datedsubsection{\textbf{Semantic Segmentation of Sequential Point Clouds
from Sparse Labels}}{Published Dec 2022}
Chenglu Wen, Weijie Huang, \textbf{Junhao Zhao}, and Cheng Wang.
\\Chinese Patent Publication No. \textit{CN115496899A}.

\section{\faTrophy\ Award}

\datedsubsection{\textbf{Second Prize, Microsoft Edge Add-ons Innovation
Challenge}}{Sep 2022}
Ranked 2nd among 148 teams. Independently developed \textit{Joy-Reading}, a
browser extension using MobileNet and TensorFlow.js to detect tables,
formulas, and other key elements for more efficient document reading.

\section{\faBriefcase\ Industry Experience}

\datedsubsection{\textbf{Software Product Manager, SND Group}, Suzhou, China}
{Jul 2023 -- Jul 2024}
Contributed to product design and development for smart agriculture and
cultural-tourism platforms supporting enterprise digital transformation.

\datedsubsection{\textbf{Data Engineer, Green Data Environmental Service
Center}}{Sep 2018 -- Mar 2019}
Suzhou, China. Collected, processed, and visualized pollution data from more than 20 cities
for a large-scale environmental monitoring system.

\section{\faMicrophone\ Selected Presentations}

\datedline{Presented at the IEEE VIS 2025 VISxGenAI workshop.}{Nov 2025}
\datedline{Attended the China 3D Vision Conference (China3DV).}{Apr 2023}

\section{\faLanguage\ Languages}

Chinese (Mandarin): native \quad
\textperiodcentered\quad English: professional working proficiency

\end{document}
